%% This section adds abstract in Finnish
\phantomsection     % this command ensures the internal links will be correct
\section*{Tiivistelmä}
\begin{otherlanguage}{finnish}
\begin{spacing}{1.5}
\universityFIN\\
\schoolFIN\\
\degreeprogrammeFIN\\
Mikäli yhteistutkinto-ohjelma, Yhteistyöyliopisto: Yliopiston nimi
\end{spacing}

\myname

\textbf{\mytitleFIN}\\
\mysubtitleFIN

\begin{spacing}{1.5}
\thesistypeFIN\\
\thesisyear\\
\pageref{myLastPage} sivua, \totalfigures~kuvaa, \totaltables~taulukkoa and XX liitettä\\
Tarkastajat: \examinerAFIN~ja \examinerBFIN\\
Avainsanat: avainsana1, avainsana2
\end{spacing}

\vspace{1em}
\begin{spacing}{1.0}
Tiivistelmätekstin riviväli on 1, ja se tasataan molempiin reunoihin, kuten muukin opinnäytetyön leipäteksti. Tiivistelmätekstin tunnistetietoineen on mahduttava yhdelle A4-arkille. Tiivistelmään ei voi sisällyttää salassa pidettäviä tietoja, vaan se on aina laadittava julkiseksi.

Tiivistelmä on itsenäinen esitys DI-työstä/pro gradu -tutkielmasta, ja sen tulee olla ymmärrettävissä sellaisenaan. Hyvässä tiivistelmässä virkkeet ovat täydellisiä, lyhyitä ja ytimekkäitä. Tekijän mielipiteet eivät näy, vaan hän kuvaa työtään kuin ulkopuolinen raportoija. Tiivistelmässä ei viitata yksityiskohtaisesti alkuperäistekstiin. Tutkimuksen keskeisimmät tulokset sisältyvät tiivistelmään.

Opiskelijat, joiden koulusivistyskieli ei ole suomi tai ruotsi ja jotka kirjoittavat opinnäytteensä englanniksi, kirjoittavat tiivistelmän ainoastaan englanniksi.
\end{spacing}
\end{otherlanguage}
\clearpage % This command will start the next section from the new page
